\chapter{复数与复变函数}

\section{复数的基本运算与初等函数}
\subsection{复数的极坐标与指数形式}
复数 \( z = x + iy \) 可表示为极坐标（指数）形式：
\[
z = |z| e^{i\arg z}
\]
其中：
\begin{itemize}
    \item \( |z| = \sqrt{x^2 + y^2} \) 为复数的模长；
    \item \( \arg z \) 为复数的辐角（主值范围：\( -\pi < \arg z \leq \pi \)）。
\end{itemize}

\subsection{对数函数}
复对数函数为多值函数，其定义为：
\[
\ln z = \ln|z| + i\arg z
\]
注：主值对数记为 \( \text{Ln}\,z = \ln|z| + i\text{Arg}\,z \)（\( \text{Arg}\,z \) 为辐角主值）。

\subsection{棣莫弗公式}
对任意正整数 \( n \)，若复数 \( z = r e^{i\theta} \)，则有棣莫弗公式：
\[
z^n = r^n e^{in\theta}
\]
展开形式：
\[
(\cos\theta + i\sin\theta)^n = \cos n\theta + i\sin n\theta
\]

\subsection{欧拉公式与三角函数}
\subsubsection{欧拉公式（核心）}
\[
e^{i\theta} = \cos\theta + i\sin\theta
\]
特殊值：当 \( \theta = \pi \) 时，\( e^{i\pi} + 1 = 0 \)（欧拉恒等式）。

\subsubsection{复变量三角函数}
由欧拉公式推导复变量三角函数定义：
\[
\cos z = \frac{e^{iz} + e^{-iz}}{2}, \quad \sin z = \frac{e^{iz} - e^{-iz}}{2i}
\]
性质：复三角函数仍满足 \( \cos^2 z + \sin^2 z = 1 \)，但无有界性（区别于实三角函数）。